\documentclass[11pt]{article}
\input{../../shared/project_style.tex}
\rhead{Basics 21 Textbook Draft}

\begin{document}
\DocTitle{Basics 21: Single-Fluid and Multifluid Descriptions}{III - Elementary plasma physics}{Species equations, collective variables, generalized Ohm's law, and the hierarchy of plasma-fluid models}

\begin{overviewbox}
\textbf{Textbook draft, version 1.0.} A plasma contains multiple species, yet many macroscopic phenomena are described by one bulk fluid. This chapter shows how species continuity and momentum equations are combined into center-of-mass, charge, and current variables, and it identifies the approximations separating multifluid, Hall-MHD, and ideal-MHD models. This chapter is designed for a reader with no prior plasma-physics training. It distinguishes exact identities from physical approximations and connects the central equations to later simulation modules.
\end{overviewbox}

\ProjectTOC

\section{Learning objectives}
After completing this note, the reader should be able to:
\begin{itemize}[leftmargin=1.6em]
\item Write species continuity and momentum equations.
\item Define mass density, bulk velocity, charge density, and current.
\item Derive the summed continuity and momentum equations.
\item Obtain the structure of generalized Ohm's law from electron momentum balance.
\item Distinguish single-fluid, two-fluid, Hall-MHD, and ideal-MHD regimes.
\item Recognize the need for pressure and heat-flux closures.
\end{itemize}

\section{Prerequisites and reading strategy}
\begin{itemize}[leftmargin=1.6em]
\item Basics 3, 5, 13, 17, and 19.
\end{itemize}
On a first reading, emphasize physical meaning and dimensional checks. On a second reading, reproduce the derivations. Attempt the computational laboratory only after the limiting cases are understood.

\section{Species fluid equations}

For each species $s$ with mass $m_s$, charge $q_s$, density $n_s$, velocity $\mathbf u_s$, and pressure tensor $\mathsf P_s$,
\begin{align}
\frac{\partial n_s}{\partial t}+\nabla\cdot(n_s\mathbf u_s)&=S_s,\\
m_sn_s\left(\frac{\partial\mathbf u_s}{\partial t}
+\mathbf u_s\cdot\nabla\mathbf u_s\right)
&=q_sn_s(\mathbf E+\mathbf u_s\times\mathbf B)
-\nabla\cdot\mathsf P_s+\mathbf R_s+\mathbf M_s.
\end{align}
$\mathbf R_s$ represents collisional momentum exchange and $\mathbf M_s$ other sources. These equations are exact moment equations only when higher moments are carried consistently.


\section{Collective variables}

Define
\begin{align}
\rho&=\sum_s m_sn_s,\\
\rho_q&=\sum_s q_sn_s,\\
\rho\mathbf u&=\sum_s m_sn_s\mathbf u_s,\\
\mathbf J&=\sum_s q_sn_s\mathbf u_s.
\end{align}
The bulk velocity is mass weighted, whereas current is charge weighted. In an electron-ion plasma with $m_e\ll m_i$, $\mathbf u$ lies close to the ion velocity even when electrons carry most of the current.


\section{Single-fluid mass conservation}

Multiply each species continuity equation by $m_s$ and sum:
\begin{equation}
\frac{\partial\rho}{\partial t}+\nabla\cdot(\rho\mathbf u)
=\sum_sm_sS_s.
\end{equation}
For ionization and recombination internal to the modeled plasma, total mass is conserved even though individual species populations change, so the right-hand side can vanish.


\section{Summed momentum equation}

Summing species momentum equations cancels internal collisional exchange:
\begin{equation}
\sum_s\mathbf R_s=0.
\end{equation}
The electromagnetic force becomes
\begin{equation}
\sum_sq_sn_s(\mathbf E+\mathbf u_s\times\mathbf B)
=\rho_q\mathbf E+\mathbf J\times\mathbf B.
\end{equation}
Under quasineutrality the electric-force term is small, giving schematically
\begin{equation}
\rho\frac{D\mathbf u}{Dt}
=\mathbf J\times\mathbf B-\nabla\cdot\mathsf P+\text{relative-drift corrections}.
\end{equation}
A strict derivation retains the stress associated with species velocities relative to $\mathbf u$.


\section{Electron momentum and generalized Ohm's law}

For singly charged electrons,
\begin{equation}
m_en_e\frac{D_e\mathbf u_e}{Dt}
=-en_e(\mathbf E+\mathbf u_e\times\mathbf B)
-\nabla\cdot\mathsf P_e+\mathbf R_e.
\end{equation}
Solving for $\mathbf E$ and writing approximately $\mathbf u_e=\mathbf u-\mathbf J/(en_e)$ yields
\begin{equation}
\mathbf E+\mathbf u\times\mathbf B
=\frac{\mathbf J\times\mathbf B}{en_e}
-\frac{\nabla\cdot\mathsf P_e}{en_e}
+\frac{\mathbf R_e}{en_e}
-\frac{m_e}{e}\frac{D_e\mathbf u_e}{Dt}.
\end{equation}
The terms are Hall, electron-pressure, resistive/collisional, and electron-inertia contributions. Their ordering determines the fluid model.


\section{Model hierarchy}

\begin{description}[leftmargin=2.3em]
\item[Multifluid:] evolves each species velocity and thermodynamics.
\item[Two-fluid:] commonly retains distinct ion and electron dynamics.
\item[Hall MHD:] keeps $\mathbf J\times\mathbf B/(en_e)$ in Ohm's law.
\item[Resistive MHD:] keeps an effective $\eta\mathbf J$ term.
\item[Ideal MHD:] assumes $\mathbf E+\mathbf u\times\mathbf B=0$.
\end{description}
These are ordered reductions, not interchangeable software options. A model is selected by comparing frequencies, length scales, collision rates, and gyro scales.


\section{Closure and loss of kinetic information}

Fluid equations contain pressure tensors and heat fluxes whose evolution introduces still higher moments. A scalar equation of state,
\begin{equation}
p=p(\rho,S),
\end{equation}
or an adiabatic law closes ideal MHD, but it removes resonant velocity-space physics, finite-Larmor-radius effects, and non-Maxwellian energetic-particle drive. Kinetic MHD restores selected kinetic responses while retaining a fluid bulk plasma.


\section{Computational laboratory}
Given electron and ion densities and velocities, compute $\rho$, $\rho_q$, $\mathbf u$, and $\mathbf J$, then reconstruct species velocities under quasineutrality. Build an ordering calculator comparing Hall, electron-inertia, pressure, and resistive terms for several characteristic scales.

\section{Summary}
\begin{itemize}[leftmargin=1.6em]
\item Multifluid equations preserve species-specific motion; single-fluid variables reorganize them into bulk and current dynamics.
\item Mass-weighted flow and charge-weighted current contain different information.
\item Generalized Ohm's law is the central discriminator among fluid plasma models.
\item Ideal MHD follows only after Hall, pressure, inertia, and resistive terms are ordered small.
\item Closure assumptions determine which kinetic effects the fluid model cannot represent.
\end{itemize}

\SymbolsAcronymsSection
\begin{symtable}
$\rho$ & mass density & $\sum_sm_sn_s$. \\
$\rho_q$ & charge density & $\sum_sq_sn_s$. \\
$\mathbf u$ & bulk velocity & Mass-weighted species velocity. \\
$\mathbf J$ & current density & Charge-weighted species flux. \\
$\mathsf P_s$ & species pressure tensor & Second central velocity moment. \\
$\mathbf R_s$ & collisional momentum exchange & Internal interspecies force density. \\
\end{symtable}

\section{Exercises}
\begin{enumerate}[leftmargin=1.8em]
\item Derive total mass conservation by summing species continuity equations.
\item Show how the species Lorentz forces combine into $\rho_q\mathbf E+\mathbf J\times\mathbf B$.
\item Derive the Hall term using $\mathbf u_e\simeq\mathbf u-\mathbf J/(en_e)$.
\item Explain why bulk velocity is generally closer to ion velocity than electron velocity.
\item Identify physics lost when a pressure tensor is replaced by scalar pressure.
\end{enumerate}

\section{References}
\begin{thebibliography}{99}
\bibitem{ref1} F. F. Chen, \emph{Introduction to Plasma Physics and Controlled Fusion}, 3rd ed., Springer (2016).
\bibitem{ref2} R. J. Goldston and P. H. Rutherford, \emph{Introduction to Plasma Physics}, CRC Press (1995).
\bibitem{ref3} P. M. Bellan, \emph{Fundamentals of Plasma Physics}, Cambridge University Press (2006).
\bibitem{ref4} J. P. Freidberg, \emph{Ideal MHD}, Cambridge University Press (2014).
\bibitem{ref5} J. P. Goedbloed, R. Keppens, and S. Poedts, \emph{Advanced Magnetohydrodynamics}, Cambridge University Press (2010).
\bibitem{ref6} H. Goedbloed and S. Poedts, \emph{Principles of Magnetohydrodynamics}, Cambridge University Press (2004).
\bibitem{ref7} D. J. Griffiths, \emph{Introduction to Electrodynamics}, 4th ed., Cambridge University Press (2017).
\bibitem{ref8} J. D. Jackson, \emph{Classical Electrodynamics}, 3rd ed., Wiley (1998).
\end{thebibliography}

\end{document}
